% Generated draft: requires mathematical and source review.
\documentclass[11pt]{article}
\usepackage[margin=25mm]{geometry}
\usepackage{fontspec}
\setmainfont{Latin Modern Roman}
\usepackage{amsmath,amssymb,mathtools,unicode-math}
\setmathfont{Latin Modern Math}
\usepackage{enumitem,needspace,xurl,hyperref}
\setlength{\parindent}{0pt}
\setlength{\parskip}{5pt}
\setlength{\emergencystretch}{3em}
\newcommand{\N}{\mathbb{N}}
\newcommand{\Z}{\mathbb{Z}}
\newcommand{\Q}{\mathbb{Q}}
\newcommand{\R}{\mathbb{R}}
\newcommand{\C}{\mathbb{C}}
\newcommand{\F}{\mathbb{F}}
\newcommand{\A}{\mathbb{A}}
\newcommand{\PP}{\mathbb{P}}
\newcommand{\E}{\mathbb{E}}
\newcommand{\eps}{\varepsilon}
\newcommand{\dd}{\,\mathrm{d}}
\begin{document}
\section*{UnsolvedMath ID: 30002586}
\textbf{Source problem code:} catalogue entry
\subsection{Definitions and mathematical statement}
The source asks to compute the noncommutative Choquet boundary and the \(C^*\)-envelope associated with a commuting operator tuple \(A=(A_1,\ldots,A_d)\). These objects require a specified operator system or operator algebra; the record does not identify that choice or the precise class of admissible tuples.

A conditional formulation in the usual bounded-operator setting is as follows. Let \(d\geq 1\) be a finite integer, let \(H\neq\{0\}\) be a complex Hilbert space, and let \(A_1,\ldots,A_d\in B(H)\) satisfy
\[
A_iA_j=A_jA_i\qquad(1\leq i,j\leq d).
\]
Put \(C_A=C^*(I_H,A_1,\ldots,A_d)\). An operator system \(\mathcal S\subseteq C_A\) is a norm-closed complex linear subspace containing \(I_H\) and closed under adjoints. Suppose that \(C^*(\mathcal S)=C_A\). One possible choice is
\[
\mathcal S_A=\operatorname{span}_{\C}\{I_H,A_1,\ldots,A_d,A_1^*,\ldots,A_d^*\},
\]
but this choice is not specified by the source. Another possible intended object is the norm-closed unital operator algebra generated by the tuple.

For a specified \(\mathcal S\), a boundary representation is an irreducible unital \(*\)-representation \(\pi:C_A\to B(K)\), with \(K\neq\{0\}\), having the following unique extension property: every unital completely positive map \(\Phi:C_A\to B(K)\) satisfying \(\Phi|_{\mathcal S}=\pi|_{\mathcal S}\) equals \(\pi\). The noncommutative Choquet boundary \(\partial_{\mathrm{nc}}\mathcal S\) consists of these representations up to unitary equivalence. Here complete positivity means that every matrix amplification of \(\Phi\) preserves positive elements.

The Shilov ideal \(J_{\mathcal S}\) is the largest closed two-sided ideal of \(C_A\) for which the quotient map is completely isometric on \(\mathcal S\): for every integer \(n\geq1\) and every \([s_{ij}]\in M_n(\mathcal S)\),
\[
\bigl\|[s_{ij}+J_{\mathcal S}]\bigr\|_{M_n(C_A/J_{\mathcal S})}
=\bigl\|[s_{ij}]\bigr\|_{M_n(C_A)}.
\]
The corresponding \(C^*\)-envelope is \(C^*_{\mathrm e}(\mathcal S)=C_A/J_{\mathcal S}\), with its induced copy of \(\mathcal S\).

\textbf{Conditional problem.} After specifying the intended operator system or operator algebra and the admissible class of commuting tuples, determine its noncommutative Choquet boundary and thereby identify its \(C^*\)-envelope. The boundedness, finite dimension of the tuple, unital convention, and operator-system choice above are explicit conventions for this conditional formulation, not additional hypotheses established by the supplied record. The source also leaves the required meaning of ``compute'' unspecified.
\subsection{Short English statement}
Determine the boundary representations and the resulting \(C^*\)-envelope associated with a commuting operator tuple. A precise target requires the intended operator system or operator algebra, the class of tuples, and the meaning of ``compute'' to be specified. The supplied record labels the problem open; that status is not independently verified here.
\subsection{Sources}
\begin{itemize}
\item Local UnsolvedMath registry, supplied complete record: permanent ID \(30002586\), source problem code \texttt{catalogue entry}, title ``Choquet Boundaries and \(C^*\)-Envelopes of Commuting Tuples.'' The record belongs to the Oberwolfach Reports Open Problems collection, gives proposed year 2014, and names no individual proposer.
\item Supplied report citation: \emph{Hilbert Modules and Complex Geometry (2015)}, Oberwolfach Reports, pp.\ 1139--1219, DOI 10.4171/owr/2014/21. \url{https://doi.org/10.4171/owr/2014/21}. The report title's supplied year 2015 differs from the record's proposed year 2014 and the year appearing in the DOI; this discrepancy is unresolved.
\item Supplied literature source: \emph{Essential normality and hyperrigidity}. \url{https://doi.org/10.1016/j.jfa.2015.03.014}.
\item Supplied literature source: \emph{Noncommutative Choquet boundary}. \url{https://arxiv.org/abs/math/0701329}.
\item Supplied literature assessment, checked \texttt{2026-08-21}: the record labels the problem open and states that general theory identifies the \(C^*\)-envelope and boundary-representation framework and computes important special cases, while a general effective computation of the Choquet boundary for the relevant commuting tuples remains open. This is supplied dated curation derived from the Oberwolfach material, not an independently verified assessment. The supplied background expressly requests independent verification of time-sensitive status and cited claims.
\end{itemize}
\end{document}
